\documentclass[11pt,a4paper]{article}
\usepackage[margin=2.2cm]{geometry}
\usepackage[T1]{fontenc}
\usepackage{lmodern}
\usepackage{xcolor}
\usepackage{booktabs}
\usepackage{tabularx}
\usepackage{enumitem}
\usepackage{listings}
\usepackage[hidelinks,colorlinks=true,linkcolor=febblue,urlcolor=febblue]{hyperref}
\definecolor{febblue}{HTML}{1C4687}
\definecolor{codebg}{HTML}{F3F4F6}
\lstset{basicstyle=\ttfamily\small, backgroundcolor=\color{codebg}, frame=single, framerule=0pt,
  framesep=4pt, xleftmargin=4pt, breaklines=true, columns=fullflexible, keepspaces=true, showstringspaces=false}
\setlength{\parskip}{4pt}
\setlist{nosep}

\title{\color{febblue}\LARGE FEBAUTO Racing: Competitor Setup (rough)\\[4pt]\large macOS, from nothing to a driving car}
\author{Formula Electric at Berkeley, Autonomous Team}
\date{September 2026 \quad\textperiodcentered\quad draft, not for distribution}

\begin{document}
\maketitle

\section*{What you are setting up}

A simulator of the RoboRacer 1:10 car runs natively on your Mac. Your code runs in a Docker
container that carries ROS~2 and the bridge to the simulator. You never install ROS on the
Mac. Foxglove shows you the car's sensors live.

\begin{lstlisting}
feb-racing/                 the repo, your clone
  feb-sim                   the launcher
  stack/                    YOUR code (gitignored)
  starter_kit/feb_driver/   the template you copy into stack/
  tracks/                   track folders
~/.feb-sim/app/mac/         the simulator app, downloaded by setup
Docker Desktop              runs the devkit container
\end{lstlisting}

The loop: the container mounts \texttt{stack/} and builds it; the bridge turns the simulator's
stream into ROS~2 topics under \texttt{/autodrive/roboracer\_1/}; your node reads the sensors
and publishes two floats, \texttt{throttle\_command} and \texttt{steering\_command}, both in
$[-1, 1]$; the bridge sends them to the simulator; the car moves.

\section{Install, once}

\begin{enumerate}
  \item \textbf{Docker Desktop.} Download from \url{https://docker.com}, install, open it, leave it running.
  \item \textbf{Foxglove.} Download from \url{https://foxglove.dev/download}, open once, sign in with a free account.
  \item \textbf{GitHub CLI.} In Terminal: \texttt{brew install gh}, then \texttt{gh auth login} and follow the prompts.
  \item \textbf{Python 3.} Already on the Mac. If Terminal says otherwise: \texttt{xcode-select --install}.
\end{enumerate}

\section{Clone and set up}

\begin{lstlisting}
git clone https://github.com/Pranman1/feb-racing
cd feb-racing
./feb-sim setup --team "Your Team"
\end{lstlisting}

Pulls the devkit image and downloads the Mac app into \texttt{\textasciitilde/.feb-sim/app/}.
A few minutes. If the app later refuses to open (``damaged'' or ``cannot be opened''), it is
unsigned; run once:

\begin{lstlisting}
xattr -dr com.apple.quarantine ~/.feb-sim/app/mac/"FEB Simulator.app"
codesign --force --deep -s - ~/.feb-sim/app/mac/"FEB Simulator.app"
\end{lstlisting}

\section{Drive manually}

\begin{lstlisting}
./feb-sim run
\end{lstlisting}

Arrow keys drive. Menu top-left: Manual/Autonomous, Track, Cars, Look, camera. Close the
window to stop. Nothing drives by itself without a stack.

\section{Drive autonomously with the starter kit}

\begin{lstlisting}
cp -r starter_kit/feb_driver stack/
./feb-sim run --stack "ros2 launch feb_driver drive.launch.py"
\end{lstlisting}

The container builds the package, the simulator opens, the car drives, Foxglove opens on its
topics. The driver is \texttt{stack/feb\_driver/feb\_driver/reactive\_driver.py}; its numbers
are in \texttt{stack/feb\_driver/config/driver.yaml}. Edit, save, run the same command again.

\section{Score and post}

\begin{lstlisting}
./feb-sim practice --stack "ros2 launch feb_driver drive.launch.py"
./feb-sim submit
\end{lstlisting}

\section{Looking inside: extra terminals and Foxglove}

\subsection*{More terminals into the devkit}

The container keeps running while the simulator is open. Open as many Terminal tabs as you
like and, in each, from the repo folder:

\begin{lstlisting}
./feb-sim shell
\end{lstlisting}

That drops you into a ROS~2 shell inside the container, with the environment already sourced.
Useful commands there:

\begin{lstlisting}
ros2 topic list                                          # everything on the graph
ros2 topic hz /autodrive/roboracer_1/lidar               # is the sensor stream alive, at what rate
ros2 topic echo /autodrive/roboracer_1/throttle_command  # what your node is sending
ros2 topic echo --once /autodrive/roboracer_1/imu
ros2 node list                                           # your node should be here
ros2 node info /driver
ros2 param list /driver                                  # live parameters
ros2 bag record -a -o /home/autodrive_devkit/src/stack/mybag   # lands in your stack/ folder on the Mac
\end{lstlisting}

\texttt{exit} leaves the shell; the container keeps running. \texttt{docker exec -it feb-devkit
bash} is the same thing without the launcher. Files you write under
\texttt{/home/autodrive\_devkit/src/stack} inside the container appear in \texttt{stack/} on the
Mac, and the other way round.

\subsection*{Foxglove}

Installed in step~2 of the install list; sign in once. After that, \texttt{./feb-sim run} opens
Foxglove connected to your car a few seconds after the simulator window. If it does not, or
you want a second window: in Foxglove choose \emph{Open connection}, \emph{Foxglove WebSocket},
address \texttt{ws://localhost:8765}.

Panels worth adding (the \emph{+} at the top right of a panel):
\begin{itemize}
  \item \textbf{3D}: add the topic \texttt{/autodrive/roboracer\_1/lidar}; set the display frame to the lidar's frame. The scan appears as points around the car.
  \item \textbf{Image}: \texttt{/autodrive/roboracer\_1/front\_camera}.
  \item \textbf{Plot}: \texttt{throttle\_command.data} and \texttt{steering\_command.data} against time; add \texttt{left\_encoder.position[0]} to watch the wheel.
  \item \textbf{Raw Messages}: any topic, to read exact values.
  \item \textbf{Topics} (left sidebar): every topic with its rate; a dash means nothing is publishing.
\end{itemize}

Save the layout (\emph{Layout}, \emph{Save}) and it comes back next time. Foxglove also opens the
bags that \texttt{practice} records under \texttt{runs/<id>/}, and any bag you record yourself:
\emph{Open local file}.

\section{When something goes wrong}

\begin{lstlisting}
docker logs feb-devkit        # build errors, your node's output
./feb-sim shell               # a ROS 2 terminal inside the container
./feb-sim stop                # kill everything and start over
\end{lstlisting}

Nothing drives: correct without \texttt{--stack}; with one, read \texttt{docker logs}.
Port 4567 in use: \texttt{./feb-sim stop}. Docker pull denied: Docker Desktop is not running.

\section{A bare-bones driver from scratch}

The smallest ROS~2 package that drives the car: it steers toward the deepest lidar beam.
It has no safety bubble and will hit walls in corners; the starter kit is the real thing.
Create \texttt{stack/my\_driver/} with these files.

\subsection*{\texttt{stack/my\_driver/package.xml}}
\begin{lstlisting}
<?xml version="1.0"?>
<package format="3">
  <name>my_driver</name>
  <version>0.0.1</version>
  <description>Bare-bones driver</description>
  <maintainer email="you@berkeley.edu">you</maintainer>
  <license>MIT</license>
  <exec_depend>rclpy</exec_depend>
  <exec_depend>sensor_msgs</exec_depend>
  <exec_depend>std_msgs</exec_depend>
  <export><build_type>ament_python</build_type></export>
</package>
\end{lstlisting}

\subsection*{\texttt{stack/my\_driver/setup.py}}
\begin{lstlisting}
from setuptools import setup
package_name = "my_driver"
setup(
    name=package_name, version="0.0.1", packages=[package_name],
    data_files=[
        ("share/ament_index/resource_index/packages", ["resource/" + package_name]),
        ("share/" + package_name, ["package.xml"]),
        ("share/" + package_name + "/launch", ["launch/drive.launch.py"]),
    ],
    install_requires=["setuptools"], zip_safe=True,
    entry_points={"console_scripts": ["driver = my_driver.driver:main"]},
)
\end{lstlisting}

\subsection*{\texttt{stack/my\_driver/setup.cfg}}
\begin{lstlisting}
[develop]
script_dir=$base/lib/my_driver
[install]
install_scripts=$base/lib/my_driver
\end{lstlisting}

\subsection*{Two empty files}
\texttt{stack/my\_driver/resource/my\_driver} and \texttt{stack/my\_driver/my\_driver/\_\_init\_\_.py}.

\subsection*{\texttt{stack/my\_driver/my\_driver/driver.py}}
\begin{lstlisting}
import math
import numpy as np
import rclpy
from rclpy.node import Node
from rclpy.qos import QoSProfile, QoSReliabilityPolicy, QoSHistoryPolicy
from sensor_msgs.msg import LaserScan
from std_msgs.msg import Float32

NS = "/autodrive/roboracer_1/"
QOS = QoSProfile(reliability=QoSReliabilityPolicy.RELIABLE,
                 history=QoSHistoryPolicy.KEEP_LAST, depth=10)
MAX_STEER = 0.5236   # rad; the command is normalised to [-1, 1]


class Driver(Node):
    def __init__(self):
        super().__init__("driver")
        self.throttle = self.create_publisher(Float32, NS + "throttle_command", QOS)
        self.steering = self.create_publisher(Float32, NS + "steering_command", QOS)
        self.create_subscription(LaserScan, NS + "lidar", self.on_scan, QOS)

    def on_scan(self, scan):
        ranges = np.asarray(scan.ranges, dtype=float)
        ranges[~np.isfinite(ranges)] = scan.range_max
        angles = scan.angle_min + np.arange(len(ranges)) * scan.angle_increment
        ahead = np.abs(angles) < math.radians(90)          # front half only
        target = angles[ahead][np.argmax(ranges[ahead])]   # bearing of the deepest beam
        steer = float(np.clip(target / MAX_STEER, -1.0, 1.0))
        speed = 0.08 if abs(steer) < 0.3 else 0.05          # 0 is a hard brake: never send it while moving
        self.steering.publish(Float32(data=steer))
        self.throttle.publish(Float32(data=speed))


def main():
    rclpy.init()
    node = Driver()
    try:
        rclpy.spin(node)
    except KeyboardInterrupt:
        pass
    node.destroy_node()
    rclpy.shutdown()
\end{lstlisting}

\subsection*{\texttt{stack/my\_driver/launch/drive.launch.py}}
\begin{lstlisting}
from launch import LaunchDescription
from launch_ros.actions import Node

def generate_launch_description():
    return LaunchDescription([
        Node(package="my_driver", executable="driver", name="driver",
             parameters=[{"use_sim_time": True}], output="screen", emulate_tty=True),
    ])
\end{lstlisting}

\subsection*{Run it}
\begin{lstlisting}
./feb-sim run --stack "ros2 launch my_driver drive.launch.py"
\end{lstlisting}

\section{What the car gives and takes}

\begin{tabularx}{\linewidth}{lX}
\toprule
Topic (under \texttt{/autodrive/roboracer\_1/}) & Meaning \\
\midrule
\texttt{lidar} & \texttt{LaserScan}, 1080 beams over 270\textdegree, 0.06 to 10~m \\
\texttt{imu} & orientation, angular velocity, linear acceleration \\
\texttt{left\_encoder}, \texttt{right\_encoder} & wheel angle in rad; radius 0.059~m \\
\texttt{front\_camera} & image \\
\texttt{throttle\_command}, \texttt{steering\_command} & \textbf{you publish these}: \texttt{Float32} in $[-1, 1]$; throttle 0 is a hard brake; steering $\pm 1 = \pm 0.5236$~rad \\
\texttt{ips}, \texttt{odom}, lap and collision counters & ground truth: training only, barred at race time \\
\bottomrule
\end{tabularx}

\smallskip
Publish as many topics of your own as you like; only the two commands reach the car. Use
message stamps or \texttt{use\_sim\_time}, never the wall clock, for speeds and integrals.
Throttle about 0.1 is roughly 2~m/s.

\end{document}
